\documentclass[11pt]{article}
\usepackage{ap-physc}
\course{AP Physics C: Electricity and Magnetism}
\headcourse{AP Physics C: E\&M}
\headpart{Unit 12 \textperiodcentered\ Question Set}
\title{Unit 12: Magnetic Fields and Electromagnetism: Question Set}
\author{Paul Wang}
\date{November 2026}
\docpart{Unit 12 \textperiodcentered\ Question Set}

\begin{document}
\maketitle

\begin{abstract}
Exam-style practice on Unit 12 (topics 12.1--12.4), which is 10--20\% of the multiple-choice section. Section~1--2 hold 13 four-choice multiple-choice questions grouped by topic; most need a short calculation or a line of reasoning, as on the exam. The last page before the key is a full Mathematical Routines free-response question (10 points). All items are original and written for practice. Use the AP equation sheet and a calculator. Allow about 26 minutes for the multiple choice (the exam's pace is about 2 minutes per question) and 24 minutes for the free response. The answer key explains every item and names its topic; a scoring guide for the free response follows it.
\end{abstract}

\scoringline{MCQ: \blank\ / \qtotal \qquad FRQ: \blank\ / 10 \qquad Target: 9+ MCQ \qquad Pace: 2 min per MCQ}

\section{Magnetic forces (12.1--12.2)}

\begin{q}{2}{12.1 Field lines}{There are no magnetic monopoles, so field lines close on themselves (running S to N inside a magnet).}
Which statement about magnetic field lines is true?
\ch{They start on north poles and end on south poles.}{They form closed loops with no beginning or end.}{They can cross where two magnets' fields meet.}{They point in the direction of the magnetic force on a moving charge.}
\end{q}
\begin{q}{4}{12.2 Force on an electron}{$\hat\imath\times\hat k = -\hat\jmath$; the electron's negative charge flips it to $+\hat\jmath$.}
An electron moves in the $+x$ direction in a magnetic field pointing in the $+z$ direction. What is the direction of the magnetic force on it?
\chh{$-y$}{$+z$}{$-x$}{$+y$}
\end{q}
\begin{q}{3}{12.2 Radius of circular motion}{$r = \dfrac{mv}{qB} = \dfrac{(1.67\times10^{-27})(2.0\times10^6)}{(1.6\times10^{-19})(0.50)} = 0.042$ m.}
A proton moving at $2.0\times10^6$ m/s enters a uniform 0.50 T field perpendicular to its velocity. What is the radius of its path? ($m_p = 1.67\times10^{-27}$ kg, $e = 1.6\times10^{-19}$ C)
\chh{$0.0023$ cm}{$2.1$ cm}{$4.2$ cm}{$8.4$ cm}
\end{q}
\begin{q}{1}{12.2 Cyclotron period}{$T = 2\pi r/v = 2\pi m/(qB)$: the radius doubles with the speed, so the time per lap does not change.}
A charged particle circles in a uniform magnetic field. If its speed is doubled, its period of revolution
\chh{stays the same}{doubles}{halves}{quadruples}
\end{q}
\begin{q}{4}{12.2 Velocity selector}{Balance: $qE = qvB$, so $v = E/B = 3000/0.020 = 1.5\times10^5$ m/s.}
Perpendicular fields $E = 3000$ V/m and $B = 0.020$ T are crossed so that their forces on a moving charge oppose. What speed passes through undeflected?
\chh{$6.7\times10^{-6}$ m/s}{$60$ m/s}{$1.5\times10^{3}$ m/s}{$1.5\times10^{5}$ m/s}
\end{q}
\begin{q}{1}{12.2 Force on a wire}{$F = ILB\sin\theta = (4.0)(0.50)(0.20)\sin 30^\circ = 0.20$ N.}
A 0.50 m wire carrying 4.0 A lies at $30^\circ$ to a uniform 0.20 T field. What is the magnitude of the force on it?
\chh{$0.20$ N}{$0.35$ N}{$0.40$ N}{$0.80$ N}
\end{q}
\begin{q}{3}{12.2 Torque on a coil}{With the plane parallel to $\vec B$, the magnetic moment (normal to the plane) is perpendicular to $\vec B$: $\tau = NIAB = (50)(2.0)(0.010)(0.30) = 0.30$ N$\cdot$m, the maximum.}
A square coil of 50 turns, 0.10 m on a side, carries 2.0 A in a 0.30 T field. The plane of the coil is parallel to the field. What is the torque on it?
\chh{$0$}{$0.0060$ N$\cdot$m}{$0.30$ N$\cdot$m}{$3.0$ N$\cdot$m}
\end{q}

\section{Fields of currents (12.3--12.4)}

\begin{q}{2}{12.3 Field of a long wire}{$B = \mu_0I/(2\pi r) = (2\times10^{-7})(20)/0.040 = 1.0\times10^{-4}$ T.}
A long straight wire carries 20 A. What is the magnetic field 4.0 cm away?
\chh{$5.0\times10^{-5}$ T}{$1.0\times10^{-4}$ T}{$3.1\times10^{-4}$ T}{$2.5\times10^{-3}$ T}
\end{q}
\begin{q}{1}{12.3 Semicircular wire}{The straight parts point along the line through P, so $d\vec l\times\hat r = 0$ for them. The semicircle gives half a full loop's field: $\tfrac12 \cdot \dfrac{\mu_0 I}{2R}$.}
A wire runs along a straight line toward a point P, around a semicircle of radius $R$ centered on P, and away from P along the same line. It carries current $I$. What is the magnetic field at P?
\ch{$\dfrac{\mu_0 I}{4R}$}{$\dfrac{\mu_0 I}{2R}$}{$\dfrac{\mu_0 I}{2\pi R}$}{$0$}
\end{q}
\begin{q}{2}{12.3 Force between wires}{$F/L = \mu_0 I_1I_2/(2\pi d) = (2\times10^{-7})(40)/0.10 = 8.0\times10^{-5}$ N/m. Same-direction currents attract.}
Two long parallel wires 0.10 m apart carry 5.0 A and 8.0 A in the same direction. What is the force per meter between them?
\chh{$8.0\times10^{-5}$ N/m, repulsive}{$8.0\times10^{-5}$ N/m, attractive}{$1.6\times10^{-4}$ N/m, attractive}{$8.0\times10^{-6}$ N/m, attractive}
\end{q}
\begin{q}{3}{12.4 Solenoid field}{$B = \mu_0 nI = (4\pi\times10^{-7})(2000)(3.0) = 7.5\times10^{-3}$ T.}
A long solenoid has 2000 turns per meter and carries 3.0 A. What is the field inside? ($\mu_0 = 4\pi\times10^{-7}$ T$\cdot$m/A)
\chh{$1.2\times10^{-3}$ T}{$3.8\times10^{-3}$ T}{$7.5\times10^{-3}$ T}{$1.5\times10^{-2}$ T}
\end{q}
\begin{q}{4}{12.4 Inside a thick wire}{$I_{enc} = I(r/R)^2 = I/4$. Ampère: $B(2\pi r) = \mu_0 I/4$ with $r = R/2$, so $B = \dfrac{\mu_0 I}{4\pi R}$.}
A long cylindrical wire of radius $R$ carries current $I$ spread uniformly over its cross-section. What is the magnetic field at $r = R/2$?
\ch{$\dfrac{\mu_0 I}{\pi R}$}{$\dfrac{\mu_0 I}{2\pi R}$}{$0$}{$\dfrac{\mu_0 I}{4\pi R}$}
\end{q}
\begin{q}{2}{12.4 Coaxial cable}{An Ampèrian circle outside encloses $I - I = 0$, and by symmetry $B$ is the same all around it, so $B = 0$.}
A coaxial cable carries current $I$ along its central wire and the same current back along its outer conductor. What is the magnetic field outside the cable?
\ch{$\dfrac{\mu_0 I}{2\pi r}$}{Zero}{$\dfrac{\mu_0 I}{\pi r}$}{$\mu_0 I$}
\end{q}

\section{Free response}

\begin{frq}{Question 1 \textperiodcentered\ Mathematical Routines \textperiodcentered\ 10 points \textperiodcentered\ 24 minutes}
In a mass spectrometer, an ion of charge $+q$ and unknown mass $m$ starts from rest and is accelerated through a potential difference $\Delta V$. It then enters a region of uniform magnetic field $B$, directed perpendicular to its velocity, travels a semicircle, and strikes a detector a distance $D$ from the point where it entered the field. Ignore gravity.

\begin{enumerate}[label=(\alph*),leftmargin=2em,itemsep=4pt]
\item Derive an expression for the ion's speed $v$ as it enters the magnetic field. \pts{2}
\item Derive an expression for the radius $r$ of the ion's path in terms of $q$, $m$, $\Delta V$, and $B$. \pts{3}
\item Derive an expression for $m$ in terms of $q$, $B$, $\Delta V$, and $D$. \pts{2}
\item Derive an expression for the time the ion spends in the magnetic field, and show that it does not depend on $\Delta V$. \pts{3}
\end{enumerate}
\end{frq}

\makekey

\subsection*{Free-response scoring guide (10 points)}
\small
Award each point only when the answer does what the model does. Points are independent unless a part says otherwise; a correct method with an answer carried forward from an earlier error earns the later points.
\begin{scoring}
(a) (2) & 1 pt: energy conservation, $q\Delta V = \tfrac12 mv^2$. 1 pt: $v = \sqrt{2q\Delta V/m}$. & Using $v = E/B$ from a velocity selector, which is not part of this device. \\
(b) (3) & 1 pt: magnetic force supplies the centripetal force, $qvB = mv^2/r$. 1 pt: $r = mv/(qB)$. 1 pt: substitutes (a): $r = \dfrac{1}{B}\sqrt{\dfrac{2m\Delta V}{q}}$. & Setting $qvB$ equal to $mv^2$ (dropping $r$) or to $mg$. \\
(c) (2) & 1 pt: $D = 2r$ (the detector is a diameter away). 1 pt: $m = \dfrac{qB^2D^2}{8\Delta V}$. & Using $D = r$. \\
(d) (3) & 1 pt: time is half a period, $t = \pi r/v$. 1 pt: with $r = mv/(qB)$, $t = \dfrac{\pi m}{qB}$. 1 pt: $\Delta V$ (hence $v$) cancels: a faster ion travels a proportionally larger semicircle. & Using $t = D/v$, a straight-line distance. \\
\end{scoring}

\end{document}
